\section*{APPENDIX}

This appendix provides proofs for the convergence results in the main text.

\subsection{Auxiliary Lemmas}

\begin{lemma}[Bounding the Client Drift]
Under Assumptions 1--3, the expected distance between the local model at step $e$ and the global model at round $t$, denoted as
\begin{equation}
D_e = \sum_{k=1}^K p_k \mathbb{E}_t\|\omega_{t,e}^k - \omega_t\|^2,
\end{equation}
is bounded by:
\begin{equation}
D_e \le \mathcal{O}(\eta_t^2 E^2) \left( C_\sigma + \|\nabla F(\omega_t)\|^2 \right),
\end{equation}
where $C_\sigma \triangleq 2\sigma_l^2 + 4\kappa^2$, and $\mathbb{E}_t[\cdot]$ denotes the expectation conditional on $\omega_t$. Consequently, the bias of the ideal gradient, $\delta_t = \mathbb{E}_t[g_t] - \nabla F(\omega_t)$, satisfies
\begin{equation}
\mathbb{E}_t\|\delta_t\|^2 \le \mathcal{O}(L^2\eta_t^2 E^2)(C_\sigma + 4M\|\nabla F(\omega_t)\|^2).
\end{equation}
\end{lemma}

\begin{proof}
From the local update rule $\omega_{t,e}^k=\omega_t-\eta_t\sum_{i=0}^{e-1}g_{t,i}^k$ and Cauchy--Schwarz,
\begin{equation}
\|\omega_{t,e}^k - \omega_t\|^2 \le \eta_t^2 e \sum_{i=0}^{e-1} \|g_{t,i}^k\|^2.
\end{equation}
Taking conditional expectation and weighted averaging over clients gives
\begin{equation}
\label{eq:bound_drift}
\sum_{k=1}^K p_k \mathbb{E}_t\|\omega_{t,e}^k - \omega_t\|^2
\le \eta_t^2 e \sum_{i=0}^{e-1} \sum_{k=1}^K p_k \mathbb{E}_t\|g_{t,i}^k\|^2.
\end{equation}
Using variance decomposition and $\|a+b\|^2 \le 2\|a\|^2+2\|b\|^2$,
\begin{equation}
\mathbb{E}_t\|g_{t,i}^k\|^2 \le 2\sigma_l^2 + 2\mathbb{E}_t\|\nabla F_k(\omega_{t,i}^k)\|^2.
\end{equation}
Adding and subtracting $\nabla F_k(\omega_t)$ and applying Assumption \ref{ass:smoothness},
\begin{equation}
\mathbb{E}_t\|g_{t,i}^k\|^2 \le 2\sigma_l^2 + 4\|\nabla F_k(\omega_t)\|^2 + 4L^2\mathbb{E}_t\|\omega_{t,i}^k - \omega_t\|^2.
\end{equation}
Summing with weights $p_k$ and invoking Assumption \ref{ass:data-heterogeneity} yields
\begin{align}
\label{eq:bound_drift_sum}
\sum_{k=1}^K p_k \mathbb{E}_t\|g_{t,i}^k\|^2
&\le \underbrace{2\sigma_l^2 + 4\kappa^2}_{C_\sigma} + 4M\|\nabla F(\omega_t)\|^2 \notag \\
&+ 4L^2 \sum_{k=1}^K p_k \mathbb{E}_t\|\omega_{t,i}^k - \omega_t\|^2.
\end{align}
Substituting (\ref{eq:bound_drift_sum}) into (\ref{eq:bound_drift}) and defining $D_e$ gives
\begin{equation}
D_e \le \eta_t^2 e \sum_{i=0}^{e-1} \left( C_\sigma + 4M\|\nabla F(\omega_t)\|^2 + 4L^2 D_i \right).
\end{equation}
For sufficiently small stepsize ($4L^2\eta_t^2 E^2 \le 1/2$), unrolling the recurrence and absorbing higher-order terms implies
\begin{equation}
D_e \le \mathcal{O}(\eta_t^2 E^2) \left( C_\sigma + \|\nabla F(\omega_t)\|^2 \right).
\end{equation}
For $\delta_t = \frac{1}{E} \sum_{e=0}^{E-1} \sum_{k=1}^K p_k (\nabla F_k(\omega_{t,e}^k) - \nabla F_k(\omega_t))$, Jensen's inequality and $L$-smoothness yield
\begin{equation}
\mathbb{E}_t\|\delta_t\|^2\!\le\!\frac{L^2}{E}\sum_{e=0}^{E-1}D_e\!\le\!\mathcal{O}\!\left(\!L^2\eta_t^2E^2\!\right)\!\left(\!C_\sigma\!+\!4M\|\nabla F(\omega_t)\|^2\!\right).
\end{equation}
\end{proof}

\begin{lemma}[Bounding the Ideal Gradient Variance]
Under Assumptions 1--3, the second moment of the ideal aggregated gradient $g_t$ is bounded by:
\begin{equation}
\mathbb{E}_t\|g_t\|^2 \le \big(1 + \mathcal{O}(L^2 \eta_t^2 E^2)\big) \cdot \Big( C_\sigma + 4M\|\nabla F(\omega_t)\|^2 \Big).
\end{equation}
\end{lemma}

\begin{proof}
By definition of $g_t$ and Jensen's inequality,
\begin{align}
\mathbb{E}_t\|g_t\|^2 &= \mathbb{E}_t\left\| \sum_{k=1}^K p_k \left( \frac{1}{E} \sum_{e=0}^{E-1} g_{t,e}^k \right) \right\|^2 \notag \\
&\le \frac{1}{E} \sum_{e=0}^{E-1} \sum_{k=1}^K p_k \mathbb{E}_t\|g_{t,e}^k\|^2.
\end{align}
Using variance decomposition,
\begin{equation}
\mathbb{E}_t\|g_{t,e}^k\|^2 \le \sigma_l^2 + \mathbb{E}_t\|\nabla F_k(\omega_{t,e}^k)\|^2.
\end{equation}
Adding and subtracting $\nabla F_k(\omega_t)$ and applying $L$-smoothness,
\begin{equation}
\mathbb{E}_t\|g_{t,e}^k\|^2 \le \sigma_l^2 + 2L^2\mathbb{E}_t\|\omega_{t,e}^k - \omega_t\|^2 + 2\|\nabla F_k(\omega_t)\|^2.
\end{equation}
Taking weighted sums and using Assumption 3,
\begin{equation}
\sum_{k=1}^K p_k \mathbb{E}_t\|g_{t,e}^k\|^2
\le \sigma_l^2 + 2\kappa^2 + 2M\|\nabla F(\omega_t)\|^2 + 2L^2 D_e.
\end{equation}
Averaging over $e$ and substituting Lemma 1,
\begin{align}
\mathbb{E}_t\|g_t\|^2 &\le \frac{1}{2}C_\sigma + 2M\|\nabla F(\omega_t)\|^2 + \frac{2L^2}{E} \sum_{e=0}^{E-1} D_e \\
&\le \frac{1}{2}C_\sigma + 2M\|\nabla F(\omega_t)\|^2 \notag \\
&+ \mathcal{O}(L^2 \eta_t^2 E^2) \left( C_\sigma + 4M\|\nabla F(\omega_t)\|^2 \right).
\end{align}
Regrouping terms gives the stated bound.
\end{proof}

\begin{lemma}[One-step Descent Lemma]
Under Assumptions 1--5, provided $\eta_t$ is sufficiently small, the expected objective value at iteration $t+1$ satisfies:
\begin{align}
\mathbb{E}_t[F(\omega_{t+1})] &\le F(\omega_t) - \mathcal{C}_0 \eta_t E \|\nabla F(\omega_t)\|^2 + \mathcal{C}_1 \eta_t^2 E^2 \notag \\
& + \mathcal{C}_2 \eta_t E B_\mu^2 + \mathcal{C}_3 \eta_t^2 E^2 B_V^2 + \tilde{R}_t,
\end{align}
where $\mathcal{C}_0, \mathcal{C}_1, \mathcal{C}_2, \mathcal{C}_3$ are positive constants, and $\tilde{R}_t = \left(\frac{8}{\eta_t E} + L\right)\|r_t\|^2$.
\end{lemma}

\begin{proof}
By Assumption 1 ($L$-smoothness),
\begin{equation}
F(\omega_{t+1}) \le F(\omega_t) + \langle \nabla F(\omega_t), \Delta_t \rangle + \frac{L}{2}\|\Delta_t\|^2.
\end{equation}
Substitute $\Delta_t = -\eta_t E(g_t + b_t) + r_t$ and decompose
\begin{align}
\mathbb{E}_t\langle \nabla F(\omega_t), \Delta_t \rangle
&= \underbrace{-\eta_t E \langle \nabla F, \mathbb{E}_t[g_t] \rangle}_{(A)}
+ \underbrace{-\eta_t E \mathbb{E}_t \langle \nabla F, b_t \rangle}_{(B)} \notag \\
&+ \underbrace{\mathbb{E}_t \langle \nabla F, r_t \rangle}_{(C)}.
\end{align}

Using $\mathbb{E}_t[g_t] = \nabla F + \delta_t$ and Young's inequality,
\begin{align}
(A) &= -\eta_t E \|\nabla F\|^2 - \eta_t E \langle \nabla F, \delta_t \rangle \notag \\
&\le -\frac{3\eta_t E}{4}\|\nabla F\|^2 + \eta_t E \|\delta_t\|^2.
\end{align}
Using $\mathbb{E}_{S_t}[b_t] = \mu_t$,
\begin{equation}
(B) \le \frac{\eta_t E}{8} \|\nabla F\|^2 + 2\eta_t E \|\mu_t\|^2 \le \frac{\eta_t E}{8} \|\nabla F\|^2 + 2\eta_t E B_\mu^2.
\end{equation}
And for $(C)$,
\begin{equation}
(C) \le \frac{\eta_t E}{16} \|\nabla F\|^2 + \frac{8}{\eta_t E}\|r_t\|^2.
\end{equation}

For the quadratic term, using $\|x+y+z\|^2 \le 3\|x\|^2 + 3\|y\|^2 + 3\|z\|^2$ and Assumption 4,
\begin{equation}
\frac{L}{2} \mathbb{E}_t\|\Delta_t\|^2 \le 2L\eta_t^2 E^2 \mathbb{E}_t\|g_t\|^2 + 2L\eta_t^2 E^2 (B_\mu^2 + B_V^2) + L\|r_t\|^2.
\end{equation}
Combining all bounds, the coefficient of $\|\nabla F(\omega_t)\|^2$ is
\begin{equation}
\text{Coeff} = -\eta_t E \left( 1 - \frac{1}{4} - \frac{1}{8} - \frac{1}{16} \right) + \text{Positive Erosion Terms.}
\end{equation}
By Lemmas 1--2, the positive erosion terms are of orders $\mathcal{O}(\eta_t^3)$ and $\mathcal{O}(\eta_t^4)$. For sufficiently small $\eta_t$ (e.g., $\eta_tEL\le c$), they are dominated by the first-order descent term, yielding a net coefficient $-\mathcal{C}_0\eta_tE$. Grouping the remaining constants gives the stated inequality.
\end{proof}

\subsection{Proof of Theorem 1}
\begin{proof}
Taking total expectation in Lemma 3 and summing over $t=0,\dots,T-1$,
\begin{align}
\mathcal{C}_0 E \sum_{t=0}^{T-1} \eta_t \mathbb{E}\|\nabla F(\omega_t)\|^2
&\le F(\omega_0) \!- \!\mathbb{E}[F(\omega_T)] \!+ \!\mathcal{C}_1 E^2 \sum_{t=0}^{T-1} \eta_t^2 \notag \\
& + \mathcal{C}_2 E B_\mu^2 \sum_{t=0}^{T-1} \eta_t + \mathcal{C}_3 E^2 B_V^2 \sum_{t=0}^{T-1} \eta_t^2 \notag \\
& + \sum_{t=0}^{T-1} \mathbb{E}[\tilde{R}_t].
\end{align}
Using $\mathbb{E}[F(\omega_T)] \ge F_{\inf}$ and $\Delta_F = F(\omega_0)-F_{\inf}$, we obtain
\begin{align}
\min_{0 \le t < T} \mathbb{E}\|\nabla F(\omega_t)\|^2
&\le \!\frac{1}{\mathcal{C}_0 E \sum \eta_t}
\Big[\sum \mathbb{E}[\tilde{R}_t]\! + \!\mathcal{C}_2 E B_\mu^2 \sum \eta_t  \notag \\
&\!+\! (\mathcal{C}_1 E^2 \!+\! \mathcal{C}_3 E^2 B_V^2)\sum \eta_t^2\!+\! \Delta_F\Big].
\end{align}
With $\eta_t = \frac{\eta_0}{\sqrt{t+1}}$,
\begin{enumerate}
\item \textbf{First-order sum:} $\sum_{t=0}^{T-1} \eta_t = \Theta(\sqrt{T})$.
\item \textbf{Second-order sum:} $\sum_{t=0}^{T-1} \eta_t^2 = \Theta(\ln T)$.
\item \textbf{Residual series:} Since $\tilde{R}_t = \left(\frac{8}{\eta_t E} + L\right)\|r_t\|^2$, and using Assumption 5 ($\|r_t\|^2 \le \mathcal{O}(t^{-2-2\rho})$):
\begin{equation}
\mathbb{E}[\tilde{R}_t] \propto \sqrt{t} \cdot t^{-2-2\rho} = t^{-1.5-2\rho}.
\end{equation}
Since $-1.5-2\rho<-1$, the $p$-series converges, i.e., $\sum_{t=0}^\infty \mathbb{E}[\tilde{R}_t] \le C_R < \infty$.
\end{enumerate}

Dividing each term on the right-hand side by $\sum \eta_t = \Theta(\sqrt{T})$:
\begin{itemize}
\item Initialization gap: $\frac{\Delta_F}{\Theta(\sqrt{T})} = \mathcal{O}\left(\frac{1}{\sqrt{T}}\right)$.
\item Variance and noise terms: $\frac{\Theta(\ln T)}{\Theta(\sqrt{T})} = \mathcal{O}\left(\frac{\ln T}{\sqrt{T}}\right)$.
\item Proxy residual: $\frac{C_R}{\Theta(\sqrt{T})} = \mathcal{O}\left(\frac{1}{\sqrt{T}}\right)$.
\item Systematic bias: $\frac{\mathcal{C}_2 E B_\mu^2 \Theta(\sqrt{T})}{\mathcal{C}_0 E \Theta(\sqrt{T})} = \mathcal{O}(B_\mu^2)$.
\end{itemize}

Combining these terms gives Theorem 1.
\end{proof}
\subsection{Proof of Corollary 1}
\begin{proof}
Taking $T\to\infty$ in Theorem 1 gives
\begin{equation}
\begin{aligned}
\lim_{T \to \infty} \Big[\, &\mathcal{O}\left( \frac{\Delta_F}{E\sqrt{T}} \right)
+ \mathcal{O}\left( \frac{(\sigma_l^2 + \kappa^2 + B_V^2) E \ln T}{\sqrt{T}} \right) \\
&+ \mathcal{O}\left( \frac{C_R}{E\sqrt{T}} \right)
+ \mathcal{O}\left( B_\mu^2 \right)\, \Big].
\end{aligned}
\end{equation}

Using standard limits,
$\lim_{T \to \infty} \frac{1}{\sqrt{T}} = 0$,
$\lim_{T \to \infty} \frac{\ln T}{\sqrt{T}} = 0$.
All terms that depend on $T$ vanish, and only the constant bias term remains:
\begin{equation}
\lim_{T \to \infty} \min_{0 \le t < T} \mathbb{E}\|\nabla F(\omega_t)\|^2 \le \mathcal{O}(B_\mu^2).
\end{equation}
thus establishing Corollary 1.
\end{proof}
%%%%%%%%%%%%%%%%%%%%%%%%%%%%%%%%%%%%%%%%%%%%%%%%%%%%%%%%%%%%%%

\bibliographystyle{IEEEtran}
\bibliography{bibfile}

%%%%%%%%%%%%%%%%%%%%%%%%%%%%%%%%%%%%%%%%%%%%%%%%%%%%%%%%%%%%%%

\begin{IEEEbiography}[{\includegraphics[width=1in,height=1.25in,clip,keepaspectratio]{fig/picture/Cuiyongsheng.jpg}}]{Yongsheng Cui}
	(Graduate Student Member, IEEE)
	is currently pursuing the Ph.D. degree in the Department of Electronic Engineering, Tsinghua University, Beijing, China.
	His research interests include deep space communications, satellite constellation design, 
	machine learning and satellite TT\&C systems.
\end{IEEEbiography}

\begin{IEEEbiography}[{\includegraphics[width=1in,height=1.25in,clip,keepaspectratio]{fig/picture/Chentaiyi.jpg}}]{Taiyi Chen}
	(Student Member, IEEE)
	received the Master degree in telecommunication engineering from the Department of Electronic Engineering, 
	Tsinghua University, Beijing, P.R. China. 
	His research interests include Wireless Communication Systems for satellite constellations, 
	channel coding, and deep space communications.
\end{IEEEbiography}

\begin{IEEEbiography}[{\includegraphics[width=1in,height=1.25in,clip,keepaspectratio]{fig/picture/Haoran Xie-eps-converted-to.pdf}}]{Haoran Xie}
	(Member, IEEE) 
    is a researcher at the Electric Power Research Institute, China Southern Power Grid.
    His research interests include application of optimization, machine learning, 
	and statistical theory. 
	He has served as the Chair for IEEE IWCMC 2022/2023/2024 Symposium and IEEE VTC 2022-Fall Symposium.
\end{IEEEbiography}

\begin{IEEEbiography}[{\includegraphics[width=1in,height=1.25in,clip,keepaspectratio]{fig/picture/YafengZhan-eps-converted-to.pdf}}]{Yafeng Zhan}
	(Member, IEEE)
	received B.S.E.E. and Ph.D.E.E. degrees from the Department of Electronic Engineering, Tsinghua University, Beijing, China, 
	in 1999 and 2004, respectively. 
	He is currently a Professor with the Beijing National Research Center for Information Science and Technology, 
	Tsinghua University. His current research interests include satellite TT\&C systems, 
	communication signal processing and deep space communications. 
    He serves as an editor of \textit{Chinese Space Science and Technology} and the \textit{Journal of Deep Space Exploration}.
\end{IEEEbiography}